\documentclass[11pt]{article}
\usepackage[a4paper,margin=23mm]{geometry}
\usepackage[T1]{fontenc}
\usepackage{lmodern,amsmath,amssymb,amsthm,mathtools,microtype,booktabs}
\usepackage[colorlinks=true,linkcolor=blue,citecolor=blue,urlcolor=blue]{hyperref}
\newcommand{\EE}{\mathbb E}
\newcommand{\tr}{\operatorname{tr}}
\newcommand{\NN}{\mathcal N}
\newcommand{\WW}{\mathcal W}
\newcommand{\LL}{\mathcal L}
\newcommand{\PP}{\mathcal P_H}
\newcommand{\QQ}{\mathcal Q}
\newcommand{\crs}{\operatorname{cr}}
\newtheorem{proposition}{Proposition}
\newtheorem{lemma}{Lemma}
\title{Oriented sources for the charged memory field\\
\large Exact transfer, whole-time bounds, and a conserved charge background}
\author{Alan Day \and Chengheng\thanks{AI research assistant (OpenAI).}}
\date{7 November 2026\quad Local research draft}
\begin{document}\maketitle
\begin{abstract}
The actual adjacent-charge transfer defect admits an exact full-Fock
loss-generator representation. We prove its weighted, compact-time
replacement by a Gaussian-number/body-loss field and give summable
ordinary-power tails. A source-dependent local normalization then
computes fixed-degree common-coupling moments, including contractions
internal to each field. Such contractions produce a conserved charge
background: even diagrams with no direct Gaussian bridge between the
fields can have a nonzero connected coefficient. We identify this
coefficient and independently verify its lowest-degree realization.
These results supply the charged input to the previously constructed
neutral kernel; the complete charged static subtraction and joint
dynamic kernel are left for the next stage.
\end{abstract}

\section{Setting and the exact adjacent-charge identity}
Use standard CAR, overlapping complex Gaussian $p$-body interactions,
the same original couplings $J$ across every charge sector, and normalized
annealed full-Fock trace at $\beta=\mu=0$. Take even $N$, $p\to\infty$,
$\lambda_N=p^2/N\to\lambda>0$, fixed $J_s$, $\hbar=1$, and string time.
Write $C=J_s^2/\lambda$, $q=e^{-4\lambda}$, and
$b_{N,m}=2^{-N}\binom Nm$. Polynomial degrees and real times are fixed
before this limit. Gaussian number $\NN$ refers to coupling chaos,
not spectral $q$-Hermite degree.

On neutral full-Fock operators define
\begin{equation}
 DX=QX-\sum_i c_i^\dagger Xc_i,\qquad Q=\sum_i c_i^\dagger c_i.
\end{equation}
For $m>0$, $h=H|_m$ and $T_mX=\sum c_i^\dagger Xc_i/m$, hence
\begin{align}
 F_m(s)&=T_m e^{isH_{m-1}}=e^{ish}-[De^{isH}]_m/m,\
 E_m(s)&=m(F_m-\alpha_m^{\NN}e^{ish})
       =[m(1-\alpha_m^{\NN})-D]e^{isH}|_m, \label{eq:exact}\
 \alpha_m&=(1-p/m)_+.
\end{align}
We set $E_0=0$; $m<p$ is inactive and $E_m=0$.
At $m=p$ the convention $0^0=1$ preserves the constant chaos.
Since $DI=0$ and $DH=pH$, $E(0)=E'(0)=0$.

\newpage
\section{Weighted replacement for the entire exponential}
Retain the finite parameter $\lambda_N$ and put
\begin{equation}
 G_N(s)=[p\NN-D-\lambda_N\NN(\NN-1)]e^{isH}.\label{eq:G}
\end{equation}
All its restrictions to $m<p$, including $m=0$, vanish.
For Gaussian matrix-valued functions use
$\|\psi\|_r=\|(\NN+1)^{r/2}\psi\|_2$, summed with $b_{N,m}$.
This is a Gaussian Sobolev norm. The preceding weighted-transfer
estimates \cite{Uniform} give, sectorwise with $C_m\le M$,
\begin{equation}
 \|H_{L/R}\psi\|_r\le A_r\sqrt M\|\psi\|_{r+1},\qquad
 \|e^{itH_L}\psi\|_r\le e^{q_r\sqrt M|t|}\|\psi\|_r,
\end{equation}
where $A_r=1+2^{r/2}$ and $q_r=r2^{\max(r-1,0)}$.

\begin{proposition}
For fixed $r,j,a,b\ge0$, $S<\infty$, and $\lambda_N$ in a positive
compact interval,
\begin{equation}
 \sup_{|s|\le S}\|H^a\partial_s^j(E-G_N)(s)H^b\|_r=O(p^{-1}).
 \label{eq:weighted}
\end{equation}
\end{proposition}
\begin{proof}
The $D$ terms cancel exactly. For $k\ge0$ and $u=p/m\le1$,
Taylor's integral remainder gives
\begin{equation}
0\le m[1-(1-u)^k]-kp+\binom k2p^2/m
 \le\binom k3p^3/m^2.\label{eq:Taylor}
\end{equation}
For $k<3$ the right side is zero. On $N/4\le m\le3N/4$ the
remaining number multiplier $R_m(k)=m(1-\alpha_m^k)-kp+
\lambda_Nk(k-1)$ obeys
\begin{equation}
 |R_m(k)|\le
 \left(\frac{4\lambda_+|m-N/2|}{N}
       +\frac{8\lambda_+^2}{3p}\right)(k+1)^3.\label{eq:R}
\end{equation}
Apply the Sobolev estimates to $H^je^{isH}$, and then the left/right
energy multipliers. The binomial identity
$\sum_m b_{N,m}(m-N/2)^2=N/4$ gives the squared error $O(p^{-2})$.
On the active complement, $|R_m(k)|\le(p+\lambda_+)(k+1)^2$;
its binomial weight is at most $2e^{-N/8}$. These estimates also
give the stated uniformity in $s$.
\end{proof}

The force interface is exact:
\begin{align}
 \QQ_\alpha f&=-f''+i\alpha\{h,f'\}+\alpha^2hfh,\
 m(\Lambda-\Lambda_{\rm OU})&=\QQ_{\alpha_m}E_m,\
 \mathsf A_m(s)&=(1-\PP)e^{-ihs/2}(\QQ_{\alpha_m}E_m)(s)e^{-ihs/2}.
\end{align}
Weighted transfer bounds $E$ and $G_N$; $\alpha_m-1=O(1/p)$ on
the broad central window. Thus $\QQ_{\alpha_m}E_m$ may be replaced
by $\QQ_1G_N$ in full-trace $L^2$. Replacing $\lambda_N$ by $\lambda$
would require the separate error $O(|\lambda_N-\lambda|)$.

\newpage
\section{An ordinary-power expansion that retains lower chaos}
Define the complete coefficients
\begin{equation}
 G_n=[p\NN-D-\lambda_N\NN(\NN-1)]H^n,\qquad
 G_N(s)=\sum_{n\ge0}\frac{(is)^n}{n!}G_n.\label{eq:rawseries}
\end{equation}
Here $G_0=G_1=0$. We do not replace this series by its highest-chaos
terms. In particular, $D$ differentiates charge-dependent contraction
coefficients.

The polar map $T=\mathcal U S_-=S_+\mathcal U$ from \cite{Transfer}
has a Gaussian-field defect bounded on chaos $k$ by
$K\sqrt{MB_m}(k+1)^2/N$, where
\begin{equation}
 B_m=\lambda_N+\lambda_N^2+
        \lambda_N(2m-N-1)^2/N.
\end{equation}
Telescoping ordinary products and using the Gaussian creation/annihilation
bound, with $L=2\sqrt M$, gives
\begin{equation}
 \|\mathcal U H_-^n-H^n\|_2
 \le\frac{K\sqrt{MB_m}}N L^{n-1}\sqrt{n!}
       \sum_{a<n}(a+1)^2.\label{eq:polar}
\end{equation}
The rank-filter estimate for $S_+-\alpha^{\NN}$ is applied to the
whole $H^n$. Multiplication by $m\le N$, followed by binomial summation,
therefore yields
\begin{equation}
 \|E_n\|_2+\|G_n\|_2\le K(n+1)^3L^n\sqrt{n!}.\label{eq:tail}
\end{equation}
The difference $E_n-G_n=R(\NN)H^n$ obeys the same type of bound by
\eqref{eq:R}. Outside the central window, polynomial factors in $N,p$
are absorbed by $e^{-N/8}$. Fixed Sobolev orders, derivatives, or
left/right energy powers only add a fixed power of $n+1$.

Consequently the force coefficients are
\begin{align}
 K_n&=G_{n+2}-\{H,G_{n+1}\}+HG_nH,\
 (\QQ_1G_N)(s)&=\sum_{n\ge0}\frac{(is)^n}{n!}K_n,\
 \|K_n\|_2&\le K'(n+3)^4L^{n+2}\sqrt{n!}.\label{eq:forcetail}
\end{align}
For every fixed $S$ the corresponding tails are bounded by sums of
$\operatorname{poly}(n)(LS)^n/\sqrt{n!}$, which converge. The same
tails hold after exact energy pinching, a Hilbert-space contraction.
They allow a fixed-degree spectral construction to be extended to
finite memory times once its static sector has been identified.

The small-time calibration is also transparent: $\mathsf A(0)$ is the
pinched $G_2$. The preceding results require this to reproduce $Y_-$;
higher coefficients must reproduce the nonzero mixed spectra in
\cite{Two}. Bounds alone do not establish those spectral identities.

\newpage
\section{Oriented marks on pure coupling-Wick sources}
For a globally neutral even source, the loss generator may be written
as a sum of symmetric local generators:
\begin{equation}
 D_i x=\tfrac12\{n_i,x\}-c_i^\dagger xc_i,
 \qquad x=\begin{pmatrix}a&b\\c&d\end{pmatrix},\qquad
 D_i x=\begin{pmatrix}0&b/2\\c/2&d-a\end{pmatrix}.\label{eq:localD}
\end{equation}
The asymmetric table is $n_ix-c_i^\dagger xc_i$; its sum agrees
on $[Q,X]=0$. The asymmetric deficit uses the raw creation count
$n_c$, whereas the symmetric one uses $(n_c+n_a)/2$.
Their equality follows from $[n_i,x]=(n_c-n_a)x$.
All fermion reordering signs must be retained before local traces.

For a $j$-slot source with no internal Gaussian pairing, let $I$ be
its incident chords. The one-mark exponents for $jp-D$, $jp-\LL$,
and $jp-\bar D$, with $\bar D=2\LL-D$, are
\begin{equation}
 A_D=j(j-1)-4\crs(I),\quad
 A_L=j^2-4\crs(I),\quad
 A_{\bar D}=j(j+1)-4\crs(I).\label{eq:pure}
\end{equation}
If two such marked sources share $b$ chords with $c$ mutual crossings,
the connected exponent $K_{12}$ is
\begin{center}
\begin{tabular}{ll}
\toprule Pair & $K_{12}$\\\midrule
$DD$ & $2b(b-1)-4c$\\
$DL$ or $LD$ & $b(b-1)-4c$\\
$LL$ & $b^2-4c$\\
$D\bar D$ & $-4c$\\
$L\bar D$ & $b(b+1)-4c$\\
$\bar D\bar D$ & $2b(b+1)-4c$\\\bottomrule
\end{tabular}
\end{center}
The double-mark weight of a complete pairing is
\begin{equation}
 C^{\#\mathrm{pairs}}q^{\crs\pi}
       [\lambda K_{12}+\lambda^2A_1A_2].\label{eq:pureweight}
\end{equation}
The one-chord $D$ deficit is identically zero, as required by $DH=pH$.
Negative mixed terms and zero unmarked traces with nonzero marked
coefficients are retained. The pure-source regression coefficients
obey $\kappa_j^D=\kappa_j^L-\lambda j$ and
$\kappa_j^{\bar D}=\kappa_j^L+\lambda j$.

These formulas are derived here using the finite-CAR and pair-Poisson
framework of \cite[Sec.~3.2--3.3]{BNR}; they are not quoted as formulas
of that paper. A nonadjacent two-slot determinant shortcut fails if
the whole source permits internal contractions. The next construction
handles exactly that case.

\newpage
\section{Whole-field sources and singleton normalization}
Fix a complete Gaussian pairing $\pi$ of $2\ell$ ordered slots.
There are $2\ell$ independent $p$-set labels; label $a$ supplies the
creator at slot $a$ and the annihilator at slot $\pi(a)$.
For a marked whole-field interval $F$ let
\begin{equation}
 i_F=\#\{\text{pairs internal to }F\},\qquad d_F=|F|-2i_F.
\end{equation}
Gaussian Wick decomposition selects degree $d_F$, including its
internal contractions. Introduce commuting nilpotents $z_F^2=0$.
The source is $\exp[z_F(pd_F-D_F)]$.
The local semigroup follows from \eqref{eq:localD}:
\begin{align}
 e^{-zD_i}c_i^{(\dagger)}&=e^{-z/2}c_i^{(\dagger)},&
 e^{-zD_i}n_i&=e^{-z}n_i,\
 e^{-zD_i}(1-n_i)&=1-e^{-z}n_i.&
\end{align}

\paragraph{Cancel singleton factors before taking the limit.}
A chord open at $F$ gives two labels, each with factor $e^{-z_F/2}$
at each of its $p$ sites. This cancels $e^{pd_Fz_F}$.
An internal chord gives the normalized singleton factors
$e^{-z_F}$ and $2-e^{-z_F}$, whose product is $1$ in the nilpotent
ring. This cancellation is exact, not an expansion leaving an $O(p)$
error. With several disjoint fields it holds in each source variable.

Let $t_a(z)$ denote the full local trace when only label $a$ occupies
a site, and $t_{ab}(z)$ the merged-site trace. Then $t_a(0)=1/2$ and
the ratio for a collision is
\begin{equation}
 R_{ab}(z)=\epsilon_{ab}\frac{t_{ab}(z)}{t_a(z)t_b(z)},\qquad
 \mathcal K_\pi(z)=\sum_{a<b}[R_{ab}(z)-1]-\ell.\label{eq:ratio}
\end{equation}
Here $\epsilon_{ab}$ is the exact fermion-grouping sign. Division is
only by singleton polynomials with nonzero constant term, never by
a possibly zero merged trace. The $-\ell$ fixes the unmarked variance
normalization. One obtains $\mathcal K_\pi(0)=-4\crs\pi$.

Pair intersections converge jointly to independent Poisson variables
of mean $\lambda$. Triple intersections vanish; fixed-degree source
coefficients are dominated by a polynomial times an exponential of
the total overlap count, whose moments are uniformly bounded.
Thus the diagram weight is $C^\ell\exp(\lambda\mathcal K_\pi)$.
For $G$ insertions shift each one-source coefficient by
$-d_F(d_F-1)$; the connected coefficients are unchanged. For two
$G$ fields this gives the explicit rule
\begin{equation}
 C^\ell q^{\crs\pi}
 \left[\lambda\kappa_{12}+
 \lambda^2(\kappa_1-d_1(d_1-1))(\kappa_2-d_2(d_2-1))\right].
 \label{eq:wholeweight}
\end{equation}
External energy slots are excluded from $d_F$. Mixed neutral marks
use their own Wick-source restrictions and the local $\LL$ table.

\newpage
\section{No direct field bridge does not mean no correlation}
Suppose two whole-field $D$ sources have no Gaussian chord directly
connecting their intervals. Their connected exponent nevertheless is
\begin{equation}
 \boxed{\ \kappa_{12}=4i_1i_2.\ }\label{eq:b0}
\end{equation}
To see this, a two-label collision can touch both fields only through
internal--internal, internal--open, or open--open chords. An open
leg scales by $e^{-z/2}$, which cancels against its source-dependent
singleton denominator; hence the latter two cases have zero mixed
coefficient. For each pair of internal chords, the four possible
label collisions have mixed coefficients $0,4,0,0$. Their sum is $4$.
This local classification is unaffected by spectator endpoints.
It also shows that a pure neutral Wick subblock has no such mixed
connected contribution in a diagram without a field bridge.

The smallest example has two $H^2$ fields, each paired internally.
It gives
\begin{equation}
 \mathcal K=-2z_1-2z_2+4z_1z_2,
 \qquad [z_1z_2]e^{\lambda\mathcal K}=4\lambda+4\lambda^2.
 \label{eq:counterexample}
\end{equation}
Subtracting only the averaged scalar leaves the positive term
$4\lambda C^2$. It is a conserved charge fluctuation, not a dynamic
correlation.

Indeed $\WW_0H^2=C_QI$ and the zero-chaos part of $G_2$ is exactly
$-DC_Q$. For the common sector variance
$C_m=w_N\binom mp\binom{N-m+p}p$, direct algebra gives
\begin{equation}
 (DC_Q)_m=m(C_m-C_{m-1})
     =C_m\frac{p(N-2m+p+1)}{N-m+1}.\label{eq:chargeexact}
\end{equation}
The formula includes $m=p$; inactive sectors are zero.
In the central window $m=N/2+xp$,
\begin{equation}
 -DC_Q=-2\lambda C(1-2x)+o(1).
\end{equation}
Under $b_{N,m}$, $X_N=(Q-N/2)/p$ tends to a centered Gaussian
with variance $1/(4\lambda)$. A uniform central linear bound and
the binomial tail give convergence of the first two moments:
\begin{equation}
 \EE\tr(-DC_Q)\to-2\lambda C,\quad
 \EE\tr(DC_Q)^2\to4C^2(\lambda+\lambda^2),\quad
 \operatorname{Var}(-DC_Q)\to4\lambda C^2.\label{eq:chargemom}
\end{equation}
This independently reproduces \eqref{eq:counterexample}. Since
$[Q,H]=0$, exact energy pinching removes this background samplewise,
including degenerate energy blocks. The general charged field still
requires a joint charge/energy regression before the neutral
no-bridge argument can be extended.

\newpage
\section{Checks, provenance, and the next calculation}
Four new scientific outputs are exclusive-create; no predecessor
driver was rerun. Their input and output hashes, dependency versions,
parameters, build logs, and final rendered pages are recorded in the
stage evidence file.
\begin{itemize}
\item Complete finite transfer: four CAR systems, 116 neutral matrix
units, 4368 Gaussian-Hermite coefficient cases, 46 independent
conditional-Gaussian OU polynomials, and 12582 exact Taylor cases.
There are 17465 matrix equalities across 468408 integer components.
The finite normalization is $C_m=2\binom mp\binom{N-m+p}p$,
distinct from the physical $J_s$ convention.
\item Oriented Wick marks: 765 geometry cases, 6120 exact coefficient
comparisons, 27 pair-site rows, 446 legal two-slot cases, and
5760 finite trace comparisons. Selected finite words include a
recorded random seed; this is not exhaustive over all index sets.
The invalid raw nonadjacent shortcut is retained as an expected
failure, with its hypothesis violation identified.
\item Whole-field sources: 6408 exact finite-CAR scalar coefficients,
with internal contractions allowed. Small systems use every index
assignment; larger cases use specified deterministic assignments.
All 294 tested no-bridge diagrams satisfy \eqref{eq:b0}; 80 have a
nonzero connected coefficient. These finite checks accompany the
local proof, not an extrapolation from a handful of large-$N$ points.
\item Charge background: 81 exact $(N,p)$ systems and 1643 sectors;
12 full binomial sums at 70 and 110 decimal digits, and 72 central
points. No charge cutoff is used in these sums. At $p=128$, the
variance relative error is still about $1.73\%$ for $\lambda=1/4$
and $8.65\%$ for $\lambda=1$. Precision agreement is not an
interval certificate or a substitute for the analytic limit.
\end{itemize}

The immediate next task is the full joint charge/energy regression
of $G_n$ and $K_n$, followed by the charged residual two-energy
measures and their mixed measures with the neutral descendants.
In particular, the known nonzero $\nu_{A_0,B_2}$ must survive.
The uniform tails \eqref{eq:tail}--\eqref{eq:forcetail} then provide
the finite-memory-time passage. The complete joint memory kernel,
the de Sitter bulk dictionary, and finite-$G$ gravitational dynamics
are not claimed here. Work remains local; this stage makes no Git
commit, push, or external issue reply.

\begin{thebibliography}{9}
\bibitem{BNR} M. Berkooz, V. Narovlansky, and H. Raj,
\emph{Complex Sachdev-Ye-Kitaev model in the double scaling limit},
\href{https://arxiv.org/abs/2006.13983}{arXiv:2006.13983},
Secs.~3.2--3.3 and Eq.~(4.3). Original fermions have anticommutator
$2$; our finite checks use standard CAR.
\bibitem{Transfer} A. Day and Chengheng,
\emph{Adjacent-charge transfer and Gaussian OU propagation},
local project stage \texttt{susskind-memory-charged-transfer} (2026).
\bibitem{Uniform} A. Day and Chengheng,
\emph{Uniform finite-memory bounds}, local project stage
\texttt{susskind-memory-uniform-bound} (2026).
\bibitem{Two} A. Day and Chengheng,
\emph{Two-component finite-memory reduction}, local project stage
\texttt{susskind-memory-two-component} (2026).
\bibitem{Neutral} A. Day and Chengheng,
\emph{A constructive neutral memory kernel}, local project stage
\texttt{susskind-memory-wick-descendants} (2026).
\end{thebibliography}
\end{document}
